\chapter{Cuplaje în grafuri bipartite}

\section{Definiții}

\begin{definition}[Graf bipartit]
  Un graf bipartit este un graf a cărui mulțime de noduri constă din două submulțimi, $V = L \cup R$, astfel încît toate muchiile duc din $L$ în $R$.
\end{definition}

\begin{definition}[Cuplaj într-un graf bipartit]
  Un cuplaj într-un graf bipartit este o colecție de muchii astfel încît orice vîrf să apară în cel mult una dintre muchii.
\end{definition}

\begin{definition}[Cuplaj maximal]
  Un cuplaj maximal este un cuplaj care nu mai poate fi extins cu nicio altă muchie din $E$.
\end{definition}

\begin{definition}[Cuplaj maxim]
  Un cuplaj maxim este un cuplaj de cardinalitate maximă.
\end{definition}

\section{Modelarea ca problemă de flux}

Similar problemelor din capitolul de fluxuri, putem modela cuplajul maxim ca problemă de flux dacă extindem graful cu următoarele elemente:

\begin{enumerate}
  \item Un nod sursă, $s$.
  \item Muchii de la $s$ la fiecare nod din $L$, de capacitate 1.
  \item Un nod destinație, $t$.
  \item Muchii de la fiecare nod din $R$ la $t$, de capacitate 1.
  \item Muchiile originale dintre $L$ și $R$ pot avea orice capacitate între 1 și $\infty$ (nu contează, căci oricum primesc doar 1 din $s$).
\end{enumerate}

Pe acest graf calculăm fluxul maxim de la $s$ la $t$. Fiecare unitate de flux va urma o cale compusă din exact trei muchii. Muchia din centru va traversa din $L$ în $R$ și o vom socoti parte din cuplaj. Cuplajul astfel obținut este maxim.

Explicația intuitivă este că forțăm fiecare nod din $L$ să fie cuplat cu cel mult un nod din $R$ și viceversa. Ele nu pot fi cuplate cu mai multe noduri, întrucît primesc cel mult o unitate de flux din $s$, respectiv transportă cel mult o unitate de flux înspre $t$.

\section{Căi alternante}

Iată și cîțiva termeni specifici cuplajelor. Îi veți regăsi în articole. Sînt doar adaptări ale denumirilor specifice fluxurilor. Fie un cuplaj $C$ (așadar o colecție de muchii disjuncte pe vîrfuri).

\begin{definition}[Nod cuplat]
  Un nod se numește cuplat dacă are o muchie adiacentă din $C$. Altfel se numește necuplat.
\end{definition}

\begin{definition}[Cale alternantă]
  O cale alternantă este o cale care traversează alternativ muchii din $C$ și din $E - C$.
\end{definition}

\begin{definition}[Drum de creștere]
  Un drum de creștere este o cale alternantă ale cărei primă și ultimă muchie fac parte din $E - C$.
\end{definition}

Observăm că drumurile de creștere au un număr impar de muchii. Ele încep și se termină cu noduri necuplate din aceeași mulțime ($L$ sau $R$).

Pentru un drum de creștere $P$, putem obține un nou cuplaj luînd diferența simetrică $C' = C \oplus P$ (diferența simetrică constă din toate elementele care sînt în $C$ sau în $P$, dar nu în ambele). Cu alte cuvinte, pentru a obține $C'$ eliminăm din $C$ toate muchiile de pe cale și le adăugăm pe celelalte. Mulțimea $C'$ este și ea un cuplaj (se verifică ușor) și conține cu o muchie în plus.

Prin fix același raționament observăm că, dacă $P$ nu este o singură cale, ci o colecție de $k$ căi disjuncte pe vîrfuri, atunci $C \oplus P$ este un cuplaj de cardinalitate $|C| + k$.

\section{Algoritmul Hopcroft-Karp}

Algoritmul lui Dinic, adaptat la grafuri bipartite, funcționează în $\bigoh(m \sqrt{n})$ și poartă numele Hopcroft-Karp. Vedeți ce simple sînt fluxurile și cuplajele? \emoji{slightly-smiling-face}

Algoritmul Hopcroft-Karp este formulat în termenii secțiunii anterioare:

\begin{enumerate}
  \item Pornește cu $C = \emptyset$.
  \item Cît timp se poate, găsește o colecție maximală $P$ de căi alternante în sensul crescător al distanțelor și adaugă colecția la $C$.
\end{enumerate}

La pasul 2 recunoaștem, de fapt, noțiunea de \textbf{blocaj de flux} din algoritmul lui Dinic. Singura diferență este că rețeaua pe niveluri are o formă particulară în grafurile bipartite. Ea constă din căi alternante.

\subsection{Analiză de complexitate}

La prima vedere, algoritmul Hopcroft-Karp are complexitatea $\bigoh(mn)$: cuplajul maxim poate avea valoarea $\min(|L|, |R|) = \bigoh(n)$, deci pot exista $\bigoh(n)$ faze. Într-o fază, găsirea unei colecții de căi alternante durează $\bigoh(m + n)$ deoarece DFS-ul nu va vizita de două ori aceeași muchie.

De fapt, putem da o limită mult mai puternică. Dacă facem căutarea căilor alternante cu BFS în rețeaua pe niveluri, atunci numărul de faze poate fi cel mult $2 \sqrt{n}$. \href{https://en.wikipedia.org/wiki/Hopcroft\%E2\%80\%93Karp_algorithm#Analysis}{Demonstrația} este mai simplă decît pare, căci refolosim ce știm de la algoritmul lui Dinic:

\begin{enumerate}
  \item Distanța de la $s$ la $t$ crește cu cel puțin 1 la fiecare fază, deci după $\sqrt{n}$ faze orice cale viitoare de la $s$ la $t$ va avea distanță mai mare decît $\sqrt{n}$.

  \item După $\sqrt{n}$ faze, ce se mai poate întîmpla pînă la găsirea cuplajului maxim? Mai putem găsi un număr de drumuri de creștere, disjuncte pe vîrfuri, de lungime cel puțin $\sqrt{n}$. De aceea, există cel mult $\sqrt{n}$ astfel de drumuri, deci algoritmul va dura cel mult încă $\sqrt{n}$ faze.
\end{enumerate}

\subsection{Detalii de implementare}

Implementarea este particularizată pentru grafuri bipartite și diferă mult de cea a algoritmului lui Dinic:

\begin{enumerate}
  \item Nu mai extinde graful cu $s$ și $t$.

  \item Nu mai dublează muchiile ca la grafurile neorientate.

  \item Nu mai construiește muchii reziduale.

  \item Așadar, ține efectiv $m$ muchii grupate în liste de adiacență doar pentru nodurile din stînga.

  \item Nodurile drepte au ca unică informație un cîmp \texttt{pair}: nodul stîng cu care sînt cuplate printr-o muchie sau 0 pentru nodurile necuplate.

  \item Nodurile stîngi au și ele cîmpul \texttt{pair} (nodul drept cu care sînt cuplate), plus lista de adiacență, plus un cîmp de distanță pentru BFS.
\end{enumerate}

Aceste informații sînt suficiente pentru a rula algoritmii BFS și DFS specifici lui Dinic / Hopcroft-Karp. Problema \hyperref[problem:school-dance]{School Dance} conține o implementare de referință.

(Menționez cu tristă veselie că unul dintre \href{https://codeforces.com/blog/entry/93824}{primele rezultate} la căutarea \textit{[Hopcroft-Karp implementation]} conține un link la cod necitibil, pe care autorul îl consideră \textit{easy to follow}. Dacă nu vă veți aminti nimic altceva din acest cerc, amintiți-vă doar să vă feriți de acest stil de codare.)

\section{Algoritmul lui Kuhn}

\href{https://cp-algorithms.com/graph/kuhn_maximum_bipartite_matching.html}{Algoritmul lui Kuhn} este mai slab din punct de vedere al complexității: $\bigoh(mn)$. Totuși, el este foarte ușor de implementat, așa că îl menționăm:

\begin{algorithm}[H]
  \caption{Algoritmul lui Kuhn}
  \begin{algorithmic}[1]
    \State inițializează un cuplaj vid
    \For{fiecare nod $u$ din mulțimea stîngă}
      \State caută cu DFS un drum de creștere care pornește din $u$
    \EndFor
  \end{algorithmic}
\end{algorithm}

\subsection{Detalii de implementare}

Problema \hyperref[problem:school-dance]{School Dance} conține o implementare completă. Algoritmul lui Kuhn stochează, pe lîngă lista de adiacență, doar două informații:

\begin{enumerate}
  \item Pentru fiecare nod drept, perechea sa din stînga (sau NIL dacă nodul nu este cuplat).
  \item Pentru fiecare nod stîng, un bit de vizitare.
\end{enumerate}

În esență, algoritmul lui Kuhn implementează metoda Ford Fulkerson cu DFS. Observăm că, în bucla \ccode{for}, nu este necesar să omitem explicit nodurile stîngi deja cuplate. Ele nu vor putea fi cuplate a doua oară, deoarece au cîmpul de vizitare deja marcat, deci DFS-ul le va ocoli.

\section{Teorema lui Kőnig}

Trebuie să acoperim, măcar în treacăt, și acest subiect, pentru că se leagă de cele discutate pînă acum. Nu vom include decît o schiță de demonstrație și nu vom face aplicații, dar sper să vă ajute să recunoașteți conceptul dacă vă veți întîlni cu el.

\begin{definition}[Acoperire cu vîrfuri]
  O acoperire cu vîrfuri a unui graf (engl. \textit{vertex cover}) este un set de vîrfuri $C$ cu proprietatea că orice muchie $e \in E$ are cel puțin un capăt în $C$.
\end{definition}

\begin{definition}[Acoperire minimă cu vîrfuri]
  O acoperire minimă cu vîrfuri (engl. \textit{MVC, minimum vertex cover}) este o acoperire cu număr minim de noduri.
\end{definition}

În grafuri generale, găsirea unei acoperiri minime este o problemă NP-hard. În grafuri bipartite, însă, putem găsi o soluție folosind \href{https://en.wikipedia.org/wiki/K\%C5\%91nig\%27s_theorem_(graph_theory)}{Teorema lui Kőnig}:

\begin{theorem}[Teorema lui Kőnig]
  În orice graf bipartit, mărimea cuplajului maxim este egală cu mărimea acoperirii minime cu vîrfuri.
\end{theorem}

Dacă vedeți o similitudine cu teorema flux maxim - tăietură minimă, aveți dreptate! Vom urmări demonstrația de pe \href{https://en.wikipedia.org/wiki/K\%C5\%91nig\%27s_theorem_(graph_theory)#Constructive_proof}{Wikipedia}, care pornește de la reducerea cuplajului maxim la problema fluxului maxim. Vom împrumuta imaginea din această demonstrație și vom detalia pașii demonstrației.

\begin{itemize}
  \item Se demonstrează ușor că orice acoperire este mai mare sau egală cu orice cuplaj. Orice cuplaj de mărime $k$ constă din $k$ muchii distincte pe noduri. Pentru a „vedea” cele $k$ muchii, orice acoperire trebuie să selecteze cel puțin $k$ noduri.

  \item De aceea, dacă putem da un exemplu de cuplaj \textbf{egal} cu o acoperire, atunci cuplajul este maxim, iar acoperirea este minimă.

  \item Dat fiind un cuplaj maxim, vom deriva o acoperire minimă, trecînd prin intermediul tăieturii minime.

  \item Ca și la flux, alegem tăietura $S-T$ în care $S$ constă din nodurile încă accesibile din $s$.

  \item Definim capacități infinite pe muchiile $A-B$, ceea ce știm că nu dăunează (nodurile din $A$ primesc oricum doar o unitate de flux din $s$).

  \item Astfel este evident că nu există muchii din $A_S$ în $B_T$: dacă ar exista, atunci tăietura $S-T$ ar avea capacitate infinită, or noi știm că ea este egală cu cuplajul maxim, care este desigur finit.

  \item (Un alt argument este că noi știm, din lecția de fluxuri, că orice muchii din $S$ în $T$ sînt saturate, ceea ce este imposibil cînd capacitatea este infinită.)

  \item Graful este orientat, deci nu există nici muchii de la $B_S$ la $A_T$.

  \item Și atunci care este capacitatea tăieturii? Este capacitatea muchiilor care traversează tăietura, adică $(s, u)$ unde $u \in A_T$ sau $(v, t)$  unde $v \in B_S$. Rezultă că și valoarea fluxului maxim este $|A_T| + |B_S|$.

  \item Apoi alegem submulțimea $A_T \cup B_S$ și demonstrăm că ea acoperă toate muchiile grafului bipartit. Într-adevăr, singura combinație de submulțimi neacoperite este $(u, v)$ cu $u \in A_S$ și $v \in B_T$, între care am demonstrat mai sus că nu există muchii.

  \item Fiind egală cu cuplajul maxim, acoperirea $A_T \cup B_S$ este minimă.
\end{itemize}
